\chapter{Software Requirements Specification of Ansieyes}

A well-defined set of requirements forms the backbone of any software system, guiding its design, implementation, and validation. In this chapter, we present the complete software requirements specification for Ansieyes, our \gls{ai}-powered issue triage system for GitHub. We outline the functional and non-functional requirements, the hardware and software dependencies, system constraints, and key assumptions that shaped our design decisions.

\section{Introduction}

Ansieyes is an intelligent issue triage system that leverages Google's Gemini 2.0 Flash \gls{llm} to automatically analyze, classify, and manage GitHub issues. The system addresses a real and growing problem in open-source and enterprise software development: the manual overhead of triaging incoming issues. By combining \gls{llm}-driven semantic understanding with traditional \gls{ml} techniques such as \gls{tfidf} cosine similarity, we built a system that can classify issues by type and severity, perform root cause analysis with code location mapping, propose solutions, detect duplicates, and defend against prompt injection attacks. The requirements documented here were gathered through a study of existing triage workflows, analysis of popular open-source repositories, and iterative prototyping.

\section{Purpose of the System}

The primary purpose of Ansieyes is to reduce the manual effort involved in triaging GitHub issues by automating the classification, analysis, and routing of incoming issues. We designed the system to serve repository maintainers, contributors, and project managers who spend significant time reading, labeling, and prioritizing issues. The system aims to provide accurate, consistent, and explainable triage decisions within seconds of an issue being opened, while also offering duplicate detection and \gls{pr} review capabilities to further streamline the development workflow.

\section{Scope of the System}

The scope of Ansieyes encompasses the following core capabilities:

\begin{itemize}
    \item Automated issue analysis and classification by type (bug, enhancement, feature request) and severity (critical, high, medium, low).
    \item Root cause analysis with mapping to specific files and code regions in the repository.
    \item Solution generation with actionable fix suggestions.
    \item Duplicate issue detection using a hybrid approach combining \gls{llm} semantic analysis and \gls{tfidf} cosine similarity.
    \item Prompt injection detection and defense through a multi-layered security module.
    \item Pull request review with automated feedback.
    \item Seamless integration with GitHub Actions for \gls{ci}/\gls{cd} pipelines.
    \item \gls{cli} tools for local and automated usage.
\end{itemize}

The system is scoped to work with public and private GitHub repositories. It does not replace human judgment but augments it by providing structured, \gls{ai}-generated analysis that maintainers can act upon.

\section{Functional Requirements}

\subsection{FR1: Issue Analysis}

The system shall accept a GitHub issue (title and body) as input and produce a structured analysis in \gls{json} format. The analysis shall include a summary of the issue, identification of affected components, and an assessment of potential impact. The system uses a two-pass Librarian-Surgeon architecture: the Librarian pass scans the full repository context to identify relevant files, and the Surgeon pass performs deep analysis on the identified subset. This approach allows Ansieyes to handle large repositories without exceeding \gls{llm} context limits.

\subsection{FR2: Issue Classification}

The system shall classify each issue into one of the following types: \textit{bug}, \textit{enhancement}, or \textit{feature request}. Additionally, it shall assign a severity level from \textit{critical}, \textit{high}, \textit{medium}, or \textit{low}. Classification results shall be validated against a Pydantic schema to ensure structural correctness before being returned.

\subsection{FR3: Root Cause Analysis}

The system shall perform root cause analysis by examining the repository's codebase in the context of the reported issue. It shall identify the most likely source files and code regions responsible for the described behavior, and provide an explanation of how the identified code relates to the issue.

\subsection{FR4: Solution Generation}

The system shall generate one or more proposed solutions for each analyzed issue. Each solution shall include a description of the fix, the files to be modified, and, where feasible, suggested code changes. Solutions are generated by the Surgeon pass after the root cause has been identified.

\subsection{FR5: Duplicate Detection}

The system shall detect whether a newly opened issue is a duplicate of an existing issue. We implement a hybrid detection strategy: first, the Gemini \gls{llm} performs semantic comparison between the new issue and existing issues; second, a \gls{tfidf} cosine similarity score is computed using scikit-learn to provide a quantitative similarity measure. An issue is flagged as a potential duplicate when both signals agree above configurable thresholds. Separate \gls{cli} commands (\texttt{ai-triage-duplicate} and \texttt{ai-triage-cosine}) are provided for each detection method.

\subsection{FR6: Prompt Injection Detection}

The system shall defend against prompt injection attacks embedded in issue titles or bodies. We employ a multi-layered security module built on the pytector library that scans all user-supplied input before it is passed to the \gls{llm}. Detected injection attempts are logged, flagged, and the malicious payload is neutralized so that the \gls{llm} processes only safe content.

\subsection{FR7: PR Review}

The system shall provide automated review of pull requests via the \texttt{ai-triage-pr} \gls{cli} command. Given a \gls{pr} identifier, the system fetches the diff, analyzes the changes using the Gemini \gls{llm}, and generates review comments covering code quality, potential bugs, and alignment with the linked issue (if any).

\subsection{FR8: GitHub Actions Integration}

The system shall integrate with GitHub Actions to enable automated triage on issue creation and \gls{pr} submission events. A reusable workflow configuration shall be provided so that repository owners can add Ansieyes to their \gls{ci}/\gls{cd} pipeline with minimal setup. The workflow shall post triage results as comments on the relevant issue or \gls{pr}.

\section{Non-Functional Requirements}

\subsection{Performance}

The system shall complete issue analysis and classification within 30 seconds for repositories containing up to 10{,}000 files. Duplicate detection shall return results within 15 seconds for repositories with up to 500 existing issues. The two-pass Librarian-Surgeon architecture ensures that \gls{llm} calls remain within acceptable latency bounds by reducing the context size fed to the Surgeon.

\subsection{Reliability}

The system shall handle \gls{api} rate limits and transient network failures gracefully by implementing retry logic with exponential backoff. Malformed \gls{llm} responses shall be caught by Pydantic validation, and the system shall retry or return a structured error rather than crashing.

\subsection{Scalability}

The system shall scale to repositories of varying sizes through its two-pass architecture. The Librarian pass uses Repomix to generate a compressed repository representation, allowing the system to process large codebases without loading all files into memory. \gls{tfidf}-based duplicate detection is computed locally using scikit-learn and scales linearly with the number of existing issues.

\subsection{Usability}

The system shall provide a clear and consistent \gls{cli} interface with four commands: \texttt{ai-triage}, \texttt{ai-triage-duplicate}, \texttt{ai-triage-cosine}, and \texttt{ai-triage-pr}. All outputs shall be structured in \gls{json} format for easy parsing by downstream tools. Configuration shall be managed via \gls{yaml} files with sensible defaults.

\subsection{Maintainability}

The codebase shall follow a modular architecture with clear separation between the \gls{llm} interaction layer, the analysis pipeline, the security module, and the GitHub integration layer. All data models shall be defined using Pydantic schemas, making it straightforward to extend or modify the system's output format. Configuration shall be externalized in \gls{yaml} files rather than hardcoded.

\section{Hardware Requirements}

\subsection{Minimum Requirements}

\begin{table}[H]
\centering
\renewcommand{\arraystretch}{1.3}
\setlength{\arrayrulewidth}{0.5mm}
\caption{\textbf{Minimum Hardware Requirements}}
\vspace{2mm}
\begin{tabular}{|p{4cm}|p{8cm}|}
\hline
\textbf{Component} & \textbf{Specification} \\
\hline
Processor & Dual-core CPU, 2.0 GHz or higher \\
\hline
RAM & 4 GB \\
\hline
Storage & 500 MB free disk space \\
\hline
Network & Stable internet connection for \gls{api} calls \\
\hline
\end{tabular}
\end{table}

\subsection{Recommended Requirements}

\begin{table}[H]
\centering
\renewcommand{\arraystretch}{1.3}
\setlength{\arrayrulewidth}{0.5mm}
\caption{\textbf{Recommended Hardware Requirements}}
\vspace{2mm}
\begin{tabular}{|p{4cm}|p{8cm}|}
\hline
\textbf{Component} & \textbf{Specification} \\
\hline
Processor & Quad-core CPU, 2.5 GHz or higher \\
\hline
RAM & 8 GB or more \\
\hline
Storage & 1 GB free disk space \\
\hline
Network & Broadband internet connection \\
\hline
\end{tabular}
\end{table}

Since the core \gls{llm} inference is offloaded to Google's Gemini \gls{api}, the local hardware requirements are modest. The primary local computation involves \gls{tfidf} vectorization and cosine similarity calculation, which are handled efficiently by scikit-learn.

\section{Software Requirements}

\begin{table}[H]
\centering
\renewcommand{\arraystretch}{1.3}
\setlength{\arrayrulewidth}{0.5mm}
\caption{\textbf{Software Requirements}}
\vspace{2mm}
\begin{tabular}{|p{4.5cm}|p{7.5cm}|}
\hline
\textbf{Component} & \textbf{Technology Used} \\
\hline
Programming Language & Python 3.11+ \\
\hline
\gls{llm} \gls{sdk} & google-genai (Google Gemini 2.0 Flash) \\
\hline
Data Validation & Pydantic \\
\hline
\gls{ml} Library & scikit-learn (\gls{tfidf} + Cosine Similarity) \\
\hline
Security Module & pytector \\
\hline
Configuration & Py\gls{yaml} \\
\hline
Repository Parsing & Repomix \\
\hline
\gls{ci}/\gls{cd} Platform & GitHub Actions \\
\hline
Version Control & Git, GitHub \\
\hline
Operating System & Linux, macOS, or Windows (with WSL) \\
\hline
\end{tabular}
\end{table}

\section{System Constraints}

\begin{enumerate}
    \item The system depends on the Google Gemini \gls{api} for all \gls{llm}-based analysis. Any downtime or quota exhaustion on the \gls{api} side will prevent the system from performing issue analysis, classification, and solution generation.
    \item The Gemini \gls{api} imposes token limits on both input and output. While the Librarian-Surgeon architecture mitigates this for large repositories, extremely large individual issues or diffs may still require truncation.
    \item GitHub \gls{api} rate limits apply when fetching issues, \gls{pr} diffs, and posting comments. Authenticated requests allow up to 5{,}000 requests per hour, which may be a constraint for repositories with very high issue volumes.
    \item The \gls{tfidf}-based duplicate detection operates on textual similarity only and may miss semantically similar issues that use very different wording. The \gls{llm}-based semantic check compensates for this limitation.
    \item The prompt injection detection module relies on known attack patterns and heuristics. Novel or highly sophisticated injection techniques may evade detection.
\end{enumerate}

\section{Assumptions}

\begin{enumerate}
    \item We assume that users have a valid Google Gemini \gls{api} key with sufficient quota for their repository's issue volume.
    \item We assume that the target GitHub repository is accessible (public, or private with appropriate token permissions) at the time of analysis.
    \item We assume that issues are written primarily in English, as the Gemini model and \gls{tfidf} vectorizer are optimized for English-language text.
    \item We assume that the repository structure follows common conventions (e.g., source code in recognizable directories) so that the Librarian pass can effectively identify relevant files.
    \item We assume a stable internet connection is available during system operation, as the system makes external \gls{api} calls to both Google and GitHub.
\end{enumerate}

\section{Summary}

In this chapter, we defined the complete software requirements specification for Ansieyes. We outlined eight functional requirements covering issue analysis, classification, root cause analysis, solution generation, duplicate detection, prompt injection defense, \gls{pr} review, and GitHub Actions integration. We specified non-functional requirements for performance, reliability, scalability, usability, and maintainability. We documented the hardware and software dependencies, identified system constraints related to \gls{api} limits and language support, and stated the assumptions under which the system is designed to operate. These requirements serve as the foundation for the system design and implementation described in the subsequent chapters.
